% EGA I Chapter0 complete3.6.2 with entire proof; French1808-1842.
% GPT-6.1 Sol / Ultra; source-linked draft, R602.
\begin{env}[3.6.2]
\label{0.3.6.2-fr}
Supongamos que $X$ es \emph{irreducible} (2.1.1)~; entonces las propiedades
siguientes son equivalentes~:
\begin{enumerate}
  \item[a)] \emph{$\mathcal{F}$ es un prehaz constante sobre $X$}~;
  \item[b)] \emph{$\mathcal{F}$ es un haz simple sobre $X$}~;
  \item[c)] \emph{$\mathcal{F}$ es un haz localmente simple sobre $X$}.
\end{enumerate}

En efecto, sea $\mathcal{F}$ un prehaz constante sobre $X$~; si $U$, $V$
son dos abiertos no vacíos de $X$, $U\cap V$ no es vacío, por lo que
$\mathcal{F}(X)\to\mathcal{F}(U)\to\mathcal{F}(U\cap V)$ y
$\mathcal{F}(X)\to\mathcal{F}(U)$ son isomorfismos, y lo mismo ocurre
con $\mathcal{F}(U)\to\mathcal{F}(U\cap V)$; análogamente,
$\mathcal{F}(V)\to\mathcal{F}(U\cap V)$ es un isomorfismo. De ello se deduce
inmediatamente que el axioma (F) de (3.1.2) se verifica, que $\mathcal{F}$ es
isomorfo a su haz asociado y, por tanto, que \emph{a)} implica \emph{b)}.

Sea ahora $(U_\alpha)$ un recubrimiento abierto de $X$ por abiertos
no vacíos y sea $\mathcal{F}$ un haz sobre $X$ tal que
$\mathcal{F}|U_\alpha$ sea simple para todo $\alpha$~; como $U_\alpha$ es
irreducible, $\mathcal{F}|U_\alpha$ es un prehaz constante en virtud de
lo anterior. Como $U_\alpha\cap U_\beta$ no es vacío,
$\mathcal{F}(U_\alpha)\to\mathcal{F}(U_\alpha\cap U_\beta)$ y
$\mathcal{F}(U_\beta)\to\mathcal{F}(U_\alpha\cap U_\beta)$ son
isomorfismos, de donde se deduce un isomorfismo canónico
$\theta_{\alpha\beta}:\mathcal{F}(U_\alpha)\to\mathcal{F}(U_\beta)$ para todo
par de índices. Pero entonces, al aplicar la condición (F) para $U=X$,
se ve que, para todo índice $\alpha_0$, $\mathcal{F}(U_{\alpha_0})$ y los
$\theta_{\alpha_0\alpha}$ son soluciones del problema universal, lo que (en
virtud de la unicidad) implica que
$\mathcal{F}(X)\to\mathcal{F}(U_{\alpha_0})$ es un isomorfismo y demuestra,
por tanto, que \emph{c)} implica \emph{a)}.
\end{env}
